% ID: SEL-D0022-OA-E02
% ESTADO: cotejada. Segunda pasada visual completa y borrador compilado revisados.
% Fuente: Selectividad/matematicas_aplicadas/reserva_a_EXAMEN 5 -ANDALUCÍA.pdf
% SHA-256: 668689642ca997cf817c9f9f6608d2835efe898e3162b050d7b22b51f225d363
% Página: 1; opción A; ejercicio original 2.

\textbf{(2.5 puntos)} Una fábrica produce entre 1000 y 6000 bombillas al día. El coste diario de producción, en euros, de $x$ bombillas viene dado por la función
\[
C(x)=9000+0.08x+\frac{2000000}{x},\qquad\text{con }1000\leq x\leq 6000.
\]
\textquestiondown{}Cuántas bombillas deberían producirse diariamente para minimizar costes? \textquestiondown{}Cuál sería dicho coste?
